\documentclass[a4paper,11pt]{article}

\usepackage[utf8]{inputenc}
\usepackage[T1]{fontenc}
\usepackage[french]{babel}
\usepackage{lmodern}
\usepackage{microtype}
\usepackage{amsmath}
\usepackage{amssymb}
\usepackage{booktabs}
\usepackage{longtable}
\usepackage{array}
\usepackage{geometry}
\usepackage{xcolor}
\usepackage{hyperref}

\geometry{margin=2.5cm}
\hypersetup{
  colorlinks=true,
  linkcolor=blue!45!black,
  urlcolor=blue!45!black,
  citecolor=blue!45!black,
  pdftitle={ASTRA-P63 Results -- Measured Ratio},
  pdfauthor={ASTRA / Representations Procedurales Adressables}
}

\setlength{\parindent}{0pt}
\setlength{\parskip}{0.65em}
\renewcommand{\arraystretch}{1.15}

\newcommand{\Code}[1]{\texttt{\detokenize{#1}}}
\newcommand{\Decision}[1]{\textbf{\texttt{\detokenize{#1}}}}

\title{\textbf{ASTRA-P63 Results}\\Measured Ratio Campaign Layer}
\author{ASTRA / Repr\'esentations Proc\'edurales Adressables}
\date{2026-05-01}

\begin{document}
\maketitle

\begin{abstract}
Ce rapport fige les r\'esultats P63-v6 du repo \Code{astra-atlas-lang}. La
trace vivante reste le rapport Markdown d'analyse P63. Ce document
\LaTeX{}/PDF est le livrable Results fig\'e pour la couche de campagnes
mesur\'ees P63. Aucun r\'esultat nouveau n'est invent\'e ici.
\end{abstract}

\section{Position dans ASTRA}

P63 appartient au chemin de calibration du ratio mesur\'e pour les
Repr\'esentations Proc\'edurales Adressables. P61 a introduit le laboratoire
virtuel avec un co\^ut proxy d\'eterministe. P62 a introduit \Code{ratio-real}
avec mesures r\'eelles par \Code{std::time::Instant} et tailles de fichiers
temporaires. P63 installe la couche durable de campagnes, de registres, de
comparaisons et de synth\`eses locales.

La question centrale reste:
\[
  R_{\mathrm{effective}} =
  \frac{\text{espace virtuel effectif, adressable et auditable}}
       {\text{co\^ut r\'eel pay\'e}}
\]

Le ratio d\'eclar\'e est affich\'e pour contexte, mais il ne pilote pas la
d\'ecision scientifique.

\section{Source et port\'ee}

La source principale est:
\begin{center}
\url{docs/analysis/ASTRA-P63-measured-ratio-calibration-analysis.md}
\end{center}

Les campagnes P63-v6 ont \'et\'e ex\'ecut\'ees localement en mode
\Code{ambitious}, avec trois campagnes de \Code{runs = 50}. Les artefacts
g\'en\'er\'es restent sous \Code{artifacts/p63/} et ne sont pas destin\'es \`a
\^etre committ\'es.

\section{R\'esultats fig\'es P63-v6}

\begin{longtable}{@{}p{0.58\linewidth}r@{}}
\toprule
M\'etrique & Valeur\\
\midrule
\Code{campaign_count} & 3\\
\Code{total_runs} & 150\\
\Code{virtual_declared_units} & 34\,320\,000\\
\Code{virtual_effective_units} & 2\,460\,000\\
\Code{total_persisted_bytes} & 777\,971\\
\Code{ratio_effective_per_byte} & 3.162072\\
\Code{gain_vs_materialized} & 352.918039\\
\Code{effective_gain_vs_materialized} & 25.296573\\
\bottomrule
\end{longtable}

\section{Cha\^ine d'adressabilit\'e virtuelle}

Les valeurs P63-v6 exposent explicitement la restriction progressive entre
l'espace d\'eclar\'e et l'espace virtuel effectivement exploitable:

\begin{longtable}{@{}lr@{}}
\toprule
Niveau & Unit\'es\\
\midrule
\Code{virtual_declared_units} & 34\,320\,000\\
\Code{virtual_reachable_units} & 3\,534\,000\\
\Code{virtual_readable_units} & 3\,054\,000\\
\Code{virtual_updatable_units} & 2\,568\,000\\
\Code{virtual_safe_units} & 2\,460\,000\\
\Code{virtual_effective_units} & 2\,460\,000\\
\bottomrule
\end{longtable}

\[
34\,320\,000 \rightarrow 3\,534\,000 \rightarrow 3\,054\,000
\rightarrow 2\,568\,000 \rightarrow 2\,460\,000 \rightarrow 2\,460\,000
\]

\section{Co\^ut r\'eel mesur\'e}

Le co\^ut r\'eel pay\'e est mesur\'e par fichiers temporaires et metadata
filesystem. La d\'ecomposition P63-v6 est:

\begin{longtable}{@{}lr@{}}
\toprule
Composant & Octets\\
\midrule
\Code{payload_file_bytes} & 136\,624\\
\Code{index_file_bytes} & 98\,378\\
\Code{journal_file_bytes} & 344\,088\\
\Code{manifest_file_bytes} & 169\\
\Code{checksum_or_audit_bytes} & 82\\
\Code{total_persisted_bytes} & 777\,971\\
\bottomrule
\end{longtable}

\Code{metadata_bytes} n'est pas s\'epar\'ement mesur\'e dans P63-v6 et reste
\Code{null} dans les rapports structur\'es.

\section{Ratios}

\begin{longtable}{@{}lr@{}}
\toprule
Ratio & Valeur\\
\midrule
\Code{ratio_declared_per_byte} & 44.114755\\
\Code{ratio_reachable_per_byte} & 4.542586\\
\Code{ratio_readable_per_byte} & 3.925596\\
\Code{ratio_updatable_per_byte} & 3.300894\\
\Code{ratio_safe_per_byte} & 3.162072\\
\Code{ratio_effective_per_byte} & 3.162072\\
\bottomrule
\end{longtable}

La m\'etrique centrale est \Code{ratio_effective_per_byte}. Le ratio d\'eclar\'e
ne suffit pas, car il ne garantit ni reachability, ni readability, ni
updatability, ni safety, ni auditabilit\'e.

\section{Baseline de mat\'erialisation}

La baseline locale suppose \Code{assumed_materialized_value_bytes = 8}. Elle
sert de comparaison transparente, pas de claim scientifique externe.

\begin{longtable}{@{}lr@{}}
\toprule
M\'etrique & Valeur\\
\midrule
\Code{assumed_materialized_value_bytes} & 8\\
\Code{estimated_materialized_bytes} & 274\,560\,000\\
\Code{gain_vs_materialized} & 352.918039\\
\Code{effective_gain_vs_materialized} & 25.296573\\
\bottomrule
\end{longtable}

\section{Vue compacte}

\begin{verbatim}
Virtual/real compression view
[declared  : ##########] 34320000 units
[effective : #---------] 2460000 units
[real cost : ##--------] 777971 bytes

ratio_effective_per_byte = 3.162072
gain_vs_materialized     = 352.918039
\end{verbatim}

\section{D\'ecisions}

\begin{longtable}{@{}ll@{}}
\toprule
Type & D\'ecision\\
\midrule
Repo / impl\'ementation & \Decision{VALIDER_IMPL_P63_MEASURED_RATIO_CAMPAIGN_LAYER}\\
Scientifique & \Decision{RECALIBRATE_P63_THRESHOLDS}\\
\bottomrule
\end{longtable}

La d\'ecision repo valide la couche d'impl\'ementation P63: exports de campagne,
registry, comparaison, campaign set summary et vue virtual/real. La d\'ecision
scientifique reste conservatrice, car \Code{allow_validate = false}, les
workloads sont internes et d\'eterministes, et les campagnes restent
mono-machine.

\section{Limites}

\begin{itemize}
  \item Une seule machine locale.
  \item Workloads internes d\'eterministes.
  \item Timings d\'ependants de la machine.
  \item Pas de datasets externes.
  \item Pas de validation scientifique P63 finale.
  \item Aucun timing goldenis\'e.
  \item Les artefacts volumineux de campagne ne sont pas committ\'es.
\end{itemize}

\section{Recommandation}

Le prochain rapport P64 devrait pr\'eparer un protocole multi-machine et des
fixtures externes, sans activer de validation scientifique avant calibration des
seuils et r\'eplication ind\'ependante.

\end{document}
